\section{Mathematische Konstruktion}

\subsection*{Leitfragen zur Recherche}
\begin{itemize}
    \item \textbf{Gitter-Methode:} Wie funktioniert die \enquote{Trapez-Methode} zur Konstruktion einer Längenanamorphose (Verzerrung eines Quadratgitters zu einem Trapezgitter)?
    \item \textbf{Strahlensatz:} Wie wendet man den Strahlensatz an, um die Abstände der Querlinien im verzerrten Raster zu berechnen?
    \item \textbf{Geometrie:} Welche geometrische Beziehung besteht zwischen dem Augpunkt (Betrachterposition), dem Bildpunkt und dem Distanzpunkt?
    \item \textbf{Projektion:} Wie lässt sich eine Anamorphose mithilfe der projektiven Geometrie beschreiben (Abbildung einer Ebene auf eine andere)?
\end{itemize}

% TODO: Hier Fließtext und ggf. Formeln/Skizzen einfügen.